\documentclass[lighthipster]{simplehipstercv}
\usepackage[utf8]{inputenc}
\usepackage[default]{raleway}
\usepackage[margin=1cm, a4paper]{geometry}

\newlength{\rightcolwidth}
\newlength{\leftcolwidth}
\setlength{\leftcolwidth}{0.23\textwidth}
\setlength{\rightcolwidth}{0.75\textwidth}

\pagestyle{empty}
\begin{document}

\thispagestyle{empty}

\section*{Inicio}

\simpleheader{headercolour}{Francisco}{Peiro}{Currículum
}{white}

\subsection*{}
\vspace{4em}

\setlength{\columnsep}{1.5cm}
\columnratio{0.23}[0.75]
\begin{paracol}{2}
\hbadness5000

\paracolbackgroundoptions

\footnotesize
{\setasidefontcolour
\flushright
\begin{center}
    \roundpic{foto.jpg}
\end{center}

\bg{cvgreen}{white}{sobre mi}\\[0.5em]

Estudiante de Ingeniería Industrial en la Universidad Nacional de Cuyo, actualmente cursando el tercer año de la carrera. Me caracterizo por ser una persona responsable, comprometida y orientada al logro de objetivos. Poseo una fuerte predisposición al aprendizaje continuo y al trabajo en equipo, así como también una buena capacidad de adaptación a nuevos entornos y desafíos. 

Además, mi participación en actividades deportivas como el rugby ha fortalecido valores como la disciplina, la constancia y el compañerismo, los cuales aplico tanto en el ámbito académico como en el profesional.

\bigskip

\bg{cvgreen}{white}{Datos Personales} \\[0.5em]
Francisco Peiro

Edad: 20 años (2006)

Nacionalidad: Argentina

\bigskip

\bg{cvgreen}{white}{Áreas de interés} \\[0.5em]

Administración ~•~ Economía ~•~ Industria ~•~ Tecnología

\bigskip

\bg{cvgreen}{white}{Intereses}\\[0.5em]

Deportes, estudios y desarrollo personal.

\bigskip

\bg{cvgreen}{white}{Herramientas}\\[0.5em]

\texttt{AutoCAD} ~/~ \texttt{Excel} ~/~ \texttt{Word} ~/~ \texttt{Octave}

\vspace{4em}

\infobubble{\faAt}{cvgreen}{white}{franpeiro15@gmail.com}
\infobubble{\faPhone}{cvgreen}{white}{2615176847}
\infobubble{\faMapMarker}{cvgreen}{white}{Mendoza, Argentina}

}

\switchcolumn

\small
\section*{Educación y Certificaciones}

\begin{tabular}{r| p{0.5\textwidth} c}
    \cvevent{2023}{Bachiller en Ciencias Económicas}{}{Colegio Corazón de María \color{cvred}}{Formación secundaria con orientación en economía y administración. Desarrollo de conocimientos en contabilidad, análisis financiero y gestión organizacional.}{blanco.jpg} \\
    
    \cvevent{2022}{Idioma Inglés}{}{Instituto Amicana \color{cvred}}{Formación en idioma inglés con enfoque en comunicación oral y escrita, permitiendo un adecuado desempeño en contextos académicos y profesionales.}{blanco.jpg}
\end{tabular}
\vspace{3em}

\begin{minipage}[t]{0.35\textwidth}
\section*{Experiencias}
\begin{tabular}{r p{0.6\textwidth} c}
    \cvdegree{2023}{Pasantías en Tecnasig SRL}{}{Área técnica-industrial \color{headerblue}}{}{blanco.jpg} \\
\end{tabular}

\vspace{0.5em}

{\small
Durante mi experiencia de pasantía en Tecnasig SRL, adquirí conocimientos prácticos en el ámbito técnico-industrial, especialmente en el área de estructuras y domos.

Participé en visitas a la refinería, lo que me permitió comprender procesos industriales reales, normas de seguridad y dinámicas de trabajo en entornos profesionales.

Asimismo, desarrollé habilidades en trabajo en equipo, responsabilidad y adaptación a contextos laborales, fortaleciendo mi perfil profesional.
}
\end{minipage}\hfill
\begin{minipage}[t]{0.3\textwidth}
\section*{División de Tiempo}
\begin{tabular}{r @{\hspace{0.5em}}l}
     \bg{skilllabelcolour}{iconcolour}{Deporte} &  \barrule{0.8}{0.5em}{cvpurple}\\
     \bg{skilllabelcolour}{iconcolour}{Estudio} & \barrule{0.7}{0.5em}{cvgreen} \\
     \bg{skilllabelcolour}{iconcolour}{Amigos y familia} & \barrule{0.6}{0.5em}{cvpurple} \\
     \bg{skilllabelcolour}{iconcolour}{Gimnasio} & \barrule{0.5}{0.5em}{cvpurple} \\
     \bg{skilllabelcolour}{iconcolour}{Lectura} & \barrule{0.4}{0.5em}{cvpurple} \\
\end{tabular}
\end{minipage}

\section*{Trayectoria}

\begin{tabular}{r| p{0.5\textwidth} c}
    \cvevent{2023}{Pasantías en Tecnasig SRL}{}{Mendoza, Argentina \color{cvred}}{Experiencia en el área técnica, participación en proyectos y visitas industriales.}{blanco.jpg} \\
\end{tabular}
\vspace{3em}

\begin{minipage}[t]{0.3\textwidth}
\section*{Certificados}
\begin{tabular}{>{\footnotesize\bfseries}r >{\footnotesize}p{0.55\textwidth}}
    2023 & Formación secundaria completa \\
    2022 & Certificación en idioma inglés
\end{tabular}
\bigskip

\section*{Idiomas}
\begin{tabular}{l | ll}
\textbf{Español} & Nativo & {\footnotesize} \\
\textbf{Inglés} & Intermedio & \pictofraction{\faCircle}{cvgreen}{2}{black!30}{2}{\tiny} \\
\end{tabular}

\end{minipage}\hfill
\begin{minipage}[t]{0.3\textwidth}
\section*{Objetivo}

{\small
Desarrollarme profesionalmente en un entorno dinámico que me permita aplicar mis conocimientos, adquirir experiencia y continuar creciendo tanto en lo personal como en lo laboral, aportando compromiso, responsabilidad y dedicación.
}
\end{minipage}

\vfill{}

\setlength{\parindent}{0pt}
\begin{minipage}[t]{\rightcolwidth}
\begin{center}\fontfamily{\sfdefault}\selectfont \color{black!70}
{\small Francisco Peiro \icon{\faEnvelopeO}{cvgreen}{} franpeiro15@gmail.com \icon{\faMapMarker}{cvgreen}{} Mendoza, Argentina \icon{\faPhone}{cvgreen}{} 2615176847}
\end{center}
\end{minipage}

\end{paracol}

\end{document}
